% SPDX-License-Identifier: GPL-3.0-or-later
% Copyright (C) 2026 Alessandro Greco
\chapter{测量网络：Tor Metrics 与 OONI}\label{cap:24-metrics}
\aggiornatoal{2026 年 9 月 23 日}{截至 2026 年 9 月 20 日的 Tor Metrics 数据；2026 年 9 月 23 日 15:00 的共识；OONI 规范 ts-023 与 ts-030}

\begin{obiettivi}
学完本章后，你将能够：
\begin{itemize}[nosep]
  \item 描述 Tor Project 发布的数据及其来源；
  \item 解释一个刻意不识别用户身份的网络如何估算用户数量，以及这些估计存在多大的不确定性；
  \item 从真实共识中提取网络统计信息（实验~15）；
  \item 解释 OONI 如何从用户视角测量 Tor 的封锁情况。
\end{itemize}
\end{obiettivi}

匿名网络给运营者提出了一个特别的问题：改进网络需要测量，但测量意味着收集的恰恰是网络本应保护的信息。Tor Project 的做法是尽量少收集数据，将结果汇总后公开，并记录计算过程的每一步，使任何人都能复现结果~\cite{metricsreproducible}。本章从数据来源出发，一直追踪到统计图表，最后介绍 OONI 如何从网络外部进行测量。

% ---------------------------------------------------------------------------
\section{数据来源}\label{sec:metrics-fonti}

公众所了解的 Tor 网络情况，几乎都来自我们已经见过的文档：
\begin{description}
  \item[共识与描述符]（\cref{cap:12-directory}）：中继列表、标志、带宽权重、版本和出口策略。
  \item[额外信息描述符（extra-info）] 中继在发布普通描述符时同时发布这些文档，其中包含历史字节传输计数，以及按国家分组的目录请求数量等汇总统计。
  \item[主动测量] OnionPerf 定期从世界各地选定的测量点通过 Tor 下载文件，记录传输时间和失败情况。
\end{description}
\textbf{CollecTor} 从网络早期起便持续归档这些文档，并提供下载。\textbf{Tor Metrics} 基于它们生成图表和数据表，附有可下载的 CSV 数据及方法说明。\textbf{Relay Search} 则支持查询单个中继。本书实验使用的共识文件（\file{fonti/collector}）即来自 CollecTor。

% ---------------------------------------------------------------------------
\section{统计匿名用户}\label{sec:metrics-utenti}

Tor 不要求用户注册，任何中继也无法知道两条连接是否属于同一个人。那么，网络是如何估计约三百万名用户的呢？方法是间接的（\cref{fig:stima}）。每个客户端都会定期下载共识，而无需识别请求者身份就可以统计目录请求数量。提供目录文档的中继每天统计成功的请求，并根据来源 IP 地址推断国家，对各国请求分别汇总。这些总数被发布在额外信息描述符中；为了避免小而精确的计数泄露更多信息，各国计数会向上取整为 8 的倍数~\spec{dir-spec}{Extra-info document format}。

\begin{figure}[htbp]
\centering
% SPDX-License-Identifier: GPL-3.0-or-later
% Copyright (C) 2026 Alessandro Greco
%% EPUB 无障碍替代文本。
\altfig{Tor Metrics 对每日平均在线客户端数量的间接估算示意，直接连接中继和通过网桥连接的客户端分别统计。提供目录服务的中继和网桥约每 24 小时统计一次成功的目录请求。extra-info 描述符中的 dirreq-v3-reqs 按国家汇总，并向上舍入为 8 的倍数。CollecTor 归档这些描述符；Tor Metrics 对有效的舍入计数减去 4，汇总观察到的请求数 r(R)，估计报告覆盖比例 frac，计算 floor(r(R)/(frac 乘以 10))。其中 10 是每名客户端平均每天请求次数的假设。覆盖比例低于 10\% 或高于 110\% 的日期会被剔除。结果估计的是平均同时在线客户端数量，不是不同自然人的人数。}
% 使用汇总目录请求间接估计平均在线客户端数，而非人数
\begin{tikzpicture}[font=\sffamily\scriptsize,node distance=0.5cm,
  box/.style={draw,rounded corners=2pt,align=center,minimum height=0.9cm,text width=2.3cm}]
  \node[box,fill=torviola!10] (r) {中继／网桥\\统计成功的\\目录请求};
  \node[box,right=of r] (x) {extra-info\\\texttt{dirreq-v3-reqs}};
  \node[box,right=of x] (c) {CollecTor\\归档数据};
  \node[box,right=of c,fill=torverde!15] (m) {Tor Metrics：\\$\lfloor r(R)/(\mathit{frac}\cdot10)\rfloor$};
  \draw[->] (r) -- (x); \draw[->] (x) -- (c); \draw[->] (c) -- (m);
  \node[below=2pt of r,font=\sffamily\tiny,text=torgrigio,align=center] {约 24 小时\\统计周期};
  \node[below=2pt of x,font=\sffamily\tiny,text=torgrigio,align=center] {按国家取整至 8 的\\倍数（再减 4）};
  \node[below=2pt of c,font=\sffamily\tiny,text=torgrigio,align=center] {$\mathit{frac}$：估计\\报告覆盖率};
  \node[below=2pt of m,font=\sffamily\tiny,text=torgrigio,align=center] {平均在线客户端\\并非独立人数};
  \node[below=0.95cm of c,font=\sffamily\tiny,text=avviso,align=center] {若 $\mathit{frac}<10\%$ 或 $\mathit{frac}>110\%$ 则剔除日期};
\end{tikzpicture}

\caption{从成功的目录请求间接估算\emph{平均在线客户端数}，而不是独立人数。目录中继和网桥分别汇总 extra-info 描述符中的请求数量；\texttt{dirreq-v3-reqs} 按国家计数并向上取整到 8 的倍数。CollecTor 归档这些数据；Tor Metrics 对有效计数减去 4 来修正舍入，再估计报告覆盖率 $\mathit{frac}$，并计算 $\lfloor r(R)/(\mathit{frac}\cdot10)\rfloor$。其中 10 是单个客户端每天平均目录请求次数的假设；若 $\mathit{frac}<10\%$ 或 $\mathit{frac}>110\%$，则排除对应日期。}
\label{fig:stima}
\end{figure}

Tor Metrics 收集这些数字，并按三个步骤处理~\cite{metricsreproducible}：
\begin{enumerate}
  \item 从各项计数中减去舍入区间的一半，以补偿向上取整造成的误差。
  \item 估算提供报告的中继覆盖全网目录请求的\textbf{比例} $\mathit{frac}$，同时考虑每台中继实际运行的时长及其处理的目录流量。
  \item 将报告中的请求数量除以该比例，估计全网总请求数，再除以 10：
  \[ r(N) = \left\lfloor \frac{r(R)}{\mathit{frac} \cdot 10} \right\rfloor. \]
\end{enumerate}
除数 10 是一种假设。始终在线的客户端每天大约请求 15 次共识，但许多客户端一天只在线一段时间，所以 Tor Project 选择 10 作为平均值。报告覆盖比例低于 10\,\% 或高于 110\,\% 的估计结果会被丢弃。

\begin{attenzione}[\enquote{用户数}究竟代表什么？]
计算结果估计的是\emph{平均在线客户端数量}，并不是独立的自然人数量。每天只运行两小时的客户端通常比全天在线的客户端产生更少的目录请求，但二者并不存在固定的比例关系；同一台机器上运行多个 Tor 程序，也会被计为多个客户端。软件行为变化，例如降低下载共识的频率，即使实际使用人数不变，也会改变估计结果。因此，这些图表更适合用来解读\emph{变化趋势}，而不是绝对人数。
\end{attenzione}

\paragraph{本书基准日期附近的数据。}2026 年 8 月 1 日至 9 月 20 日，直接连接 Tor 的估计日用户数在 280 万到 350 万之间波动，报告覆盖比例为 57\,\%--73\,\%。意大利的日均估计值约为 60\,000，但每日结果的预期区间很宽。例如 9 月 20 日，区间约为 34\,000--111\,000。在这样的范围内，数字变化本身不足以证明发生了特定事件。同一天，按传输技术估计的网桥用户数分别约为：obfs4 98\,000、Snowflake 60\,000、WebTunnel 35\,000。

\subsection{发现网络审查}

Metrics 使用相同的时间序列标记可能的审查事件。对于用户数量足够多的国家，它将该国的周环比变化与其他国家的典型变化比较，并突出显示超出预期区间的日期~\cite{metricsreproducible}。直接连接的用户数骤降，同时网桥用户数上升，是遭到封锁的典型信号。但这只是一条线索，不能视为证据。确认情况需要来自该国境内的测量结果，而这正是 OONI 的工作。

% ---------------------------------------------------------------------------
\section{性能测量}\label{sec:metrics-prestazioni}

OnionPerf 从固定位置测量通过 Tor 下载不同大小文件所需的时间、失败情况，以及电路延迟：从发送 HTTP 请求到收到响应头之间的时间~\cite{metricsreproducible}。2026 年 9 月 20 日，从德国测量点访问公共服务器的延迟中位数约为 110--140\,ms，美国测量点约为 280\,ms，香港测量点则超过 500\,ms。测量点与欧洲、北美主要 Tor 网络资源之间的地理距离具有重要影响。部分测量点使用 \cref{cap:31-futuro} 介绍的多路径机制 Conflux，从而可以在相近条件下进行比较。

% ---------------------------------------------------------------------------
\section{实验：读取真实共识}\label{sec:lab-metrics}

\begin{laboratorio}[15]{从共识中推导网络统计信息}
\textbf{目标。}使用新模块 \file{labs/zh-Hans/代码/torlab/metrics.py}，从 CollecTor 的真实共识中提取关键网络统计指标。\\
\textbf{前置条件。}实验~6（\file{directory.py}）和实验~7（\file{path\_selection.py}）；无需启动私有网络。\\
\textbf{预计时间。}1 小时。
\end{laboratorio}

该模块新增 \texttt{summary} 函数，复用实验~6 的共识解析器，同时提供两个用户估计函数。当数值相同时，统计结果也会按名称排序。否则，Python 集合的迭代顺序可能导致每次执行的结果顺序不同。要复现 2026 年 9 月 23 日 15:00 的历史共识结果，请先在仓库根目录运行 \texttt{make -C labs/zh-Hans data} 准备数据（下载时需要联网）。应用于该共识（输出有节选）：
\begin{lstlisting}[language=bash,numbers=none]
$ PYTHONPATH=labs/zh-Hans/代码 python3 -m torlab.metrics labs/zh-Hans/数据/collector/consensus-microdesc-2026-09-23-15
relay: 9358
flag: {'Running': 9358, 'Valid': 9358, 'Fast': 8961, 'Stable': 8199,
       'V2Dir': 7674, 'Guard': 5306, 'HSDir': 4050, 'Exit': 3177,
       'BadExit': 32, 'MiddleOnly': 28, 'StaleDesc': 21, 'Authority': 10}
total_bandwidth: 198927900
exit_bandwidth_share: 0.3308947161257923
guard_bandwidth_share: 0.7010271560701139
versions: [('Tor 0.4.9.12', 6027), ('Tor 0.4.9.11', 2520), ...]
networks16: [('64.65', 771), ('192.42', 636), ('204.8', 598), ...]
\end{lstlisting}

可以观察到以下几点：
\begin{itemize}
  \item 共识列出的 9\,358 台中继，与 Tor Metrics 图表相符：9 月 20 日，图表显示 9\,416 台中继和 2\,333 台网桥。
  \item 大约三分之一的中继拥有 \texttt{Exit} 标志，约占所测量带宽的三分之一；拥有 Guard 标志的中继约占 70\,\% 的带宽。这两个比例可以重叠，因为一台中继可能同时拥有两个标志。\cref{sec:percorso-pesi} 中介绍的带宽权重用于处理这些不平衡。
  \item \texttt{BadExit}（32 台中继）和 \texttt{MiddleOnly}（28 台）反映目录权威节点对可疑或配置不当中继的判断。
  \item 有十台中继带有 \texttt{Authority} 标志，但只有九台权威节点参与共识投票（\cref{sec:arch-ruoli}）。第十台是网桥权威节点 \texttt{Serge}，虽带有该标志，却不参与投票。
  \item 多数中继运行较新的 0.4.9 系列版本。
  \item 部分 /16 网络包含数百台中继，通常是大型出口运营者。\cref{sec:percorso-vincoli} 中的每个 /16 网络最多选择一台中继的规则，可以减少电路对单一网络的依赖。
\end{itemize}
共识带宽权重的总和代表目录权威节点测量的带宽单位，而不是实际传输的流量。要估算真实流量，Metrics 使用额外信息描述符中的历史字节计数。

% ---------------------------------------------------------------------------
\section{OONI}\label{sec:metrics-ooni}

\emph{Open Observatory of Network Interference}（OONI，开放网络干扰观测站）从用户视角测量互联网审查。世界各地的志愿者安装 OONI Probe，定期执行测试并公开结果。其中有些测试专门针对 Tor~\cite{oonitor}：
\begin{description}
  \item[\texttt{tor}] 检查选定目录权威节点和网桥的可达性。测试不仅仅尝试建立 TCP 连接，还执行协议特有的操作，例如 obfs4 握手。
  \item[\texttt{torsf}] 尝试通过 Snowflake 启动 tor，记录引导是否成功及耗时。
\end{description}
这些测试规范包含防止泄露私人网桥地址的保护措施：如果将这些地址写入公开结果，就等于直接交给审查者。Metrics 和 OONI 互为补充。Metrics 可能显示某个国家的用户数量下降；OONI 则能调查该国境内是否真的无法连接中继或网桥，并帮助识别封锁方式。

\begin{approfondimento}[测量不能造成伤害]
中继发布的每项统计数据，都是在实用性与隐私之间进行权衡。细节越多，数据越有用，泄露信息的风险也越大。因此，请求计数会进行舍入，并按日期和国家汇总。网络研究也应遵循同样原则（\cref{sec:rp-perche}）：只收集必要的最少数据，在发布前进行汇总，不保留不必要的资料。
\end{approfondimento}

% ---------------------------------------------------------------------------
\section*{本章小结}
\begin{itemize}
  \item 网络数据来自 CollecTor 归档的共识、描述符与额外信息文档，以及 OnionPerf 主动测量。Tor Metrics 使用公开、可复现的方法分析它们。
  \item 用户数量依据目录请求估算：先用报告数量除以覆盖比例，再除以 10。结果表示平均在线客户端数，尤其适合研究变化趋势。
  \item 同样的数据还可能提示审查事件，但需要当地测量来确认。
  \item 真实共识可以提供中继数量、标志、带宽占比、版本和网络集中程度。
  \item OONI 从不同国家境内测量权威节点与网桥的可达性，以及 Snowflake 的运行情况。
\end{itemize}

% ---------------------------------------------------------------------------
\section*{练习}

\begin{esercizio}[1]\label{es:24-dieci}
如果所有客户端的共识下载频率都减少一半，而实际使用人数保持不变，用户数量估计会如何变化？解读图表的人可能怎样察觉这一变化？
\end{esercizio}

\begin{esercizio}[1]\label{es:24-frazione}
某一天，提供报告的中继记录了两千万次请求，估计报告覆盖比例为 60\,\%。Metrics 会估算出多少用户？如果覆盖比例是 8\,\%，又会怎样处理？
\end{esercizio}

\begin{esercizio}[2]\label{es:24-flag}
扩展 \texttt{summary} 函数，计算同时带有 \texttt{Guard} 和 \texttt{Exit} 标志的中继所占的带宽比例，并估算电路出口属于中继数量最多的五个 /16 网络之一的加权概率。
\end{esercizio}

\begin{esercizio}[2]\label{es:24-censura}
选择一个发生过明确记录的 Tor 封锁事件的国家，从 Metrics 下载直接连接用户和网桥用户的时间序列。描述封锁发生时两个序列如何变化。
\end{esercizio}

\begin{esercizio}[3]\label{es:24-privacy}
各国请求计数会向上舍入为 8 的倍数。描述一种因此更难实施的攻击，并解释为什么这一措施无法抵御控制大量中继的攻击者。
\end{esercizio}

% ---------------------------------------------------------------------------
\section*{延伸阅读}
\begin{itemize}
  \item Tor Project 的 \titolo{Reproducible Metrics} 页面~\cite{metricsreproducible}，介绍各类图表背后的方法。
  \item OONI 测试规范~\cite{oonitor}，以及 OONI Explorer：\url{https://explorer.ooni.org/}。
  \item 额外信息描述符格式：\spec{dir-spec}{Extra-info document format}。
\end{itemize}
